\documentclass[12pt,a4paper]{article}
\usepackage{ucebnice}
\usepackage{polynom}

\title{Kvadratické funkce}
\author{Ondřej Škubala}
\date{\today}

\begin{document}

\maketitle
\newpage

\section{Kvadratická funkce}

\begin{quotebox}
{„Parabola je zrcadlem, ve kterém se matematika dívá sama na sebe.“}
{René Descartes, \textit{La Géométrie (1637)}}
\end{quotebox}

\begin{historicalbox}
\textbf{Parabola a její příběh napříč staletími}

Historie paraboly je takřka stejně fascinující jako samotný tvar, který opisuje.  
Už ve 4.~století př.~n.~l. se řečtí matematici zabývali otázkou, jaké křivky vznikají při průniku kužele s rovinou.  
Byl to žák Platóna, \textbf{Menaichmos}, kdo poprvé rozpoznal zvláštní vlastnosti této křivky a zavedl její popis pomocí geometrických vztahů.  
Zajímalo ho, zda lze s její pomocí vyřešit slavný problém \textbf{zdvojení krychle} — tedy sestrojit hranu nové krychle, jejíž objem by byl přesně dvojnásobkem té původní.  
Ačkoliv se to nakonec ukázalo být nemožné pouze s kružítkem a pravítkem, právě tato snaha vedla ke zrodu nového druhu křivek – kuželoseček.

O několik století později se problematikou zabýval \textbf{Apollónios z Pergy}, často označovaný jako „Velký geometer“.  
Ve svém díle \textit{Kuželosečky} popsal detailně nejen parabolu, ale i elipsu a hyperbolu.  
Jeho texty jsou pro dnešního čtenáře překvapivě moderní: systematicky rozlišuje jednotlivé případy podle polohy roviny a kužele, zavádí terminologii, kterou používáme dodnes, a naznačuje i určitou symetrii mezi těmito třemi typy křivek.  
Jeho pojetí paraboly se však stále opíralo čistě o geometrické vlastnosti, nikoliv o algebraické vztahy.

\smallskip
Teprve o mnoho staletí později, v období renesance, se parabola objevuje znovu — tentokrát už jako \textbf{fyzikální realita}.  
Italský učenec \textbf{Galileo Galilei} zjistil, že dráha tělesa, které se pohybuje ve vakuu pouze vlivem gravitace, je právě parabolická.  
Tím poprvé spojil \textbf{geometrii s přírodními zákony}, čímž se parabola stala nejen matematickým objektem, ale i nástrojem k pochopení světa kolem nás.  
Od té doby ji nacházíme všude: v pohybu projektilů, světelných paprsků i proudění tekutin.

\smallskip
Zásadní obrat nastal s příchodem \textbf{Reného Descarta} a jeho knihy \textit{La Géométrie} (1637), která položila základy \textbf{analytické geometrie}.  
Descartes ukázal, že každou parabolu lze popsat pomocí \textbf{kvadratické rovnice}, a tím spojil dávnou řeckou intuici s novou algebraickou metodou.  
Křivka, kterou Apollónios kreslil kružítkem, se tak stala obrazem rovnice \(y = ax^2 + bx + c\).  
Tím se zrodilo něco mimořádného: jazyk, v němž geometrie a algebra začaly vzájemně hovořit.

\smallskip
V následujících stoletích se parabola stala nejen matematickým symbolem, ale i \textbf{architektonickým a technickým prvkem}.  
Její vlastnost soustředit všechny paprsky přicházející rovnoběžně s osou do jednoho ohniska využívají reflektory, satelitní antény i parabolické mikrofony.  
V moderní době se navíc objevuje ve zcela novém kontextu: paraboly popisují zisky firem, trajektorie planet, a v informatice se dokonce používají v algoritmech pro optimalizaci a strojové učení.

\smallskip
Z malého geometrického experimentu starých Řeků se tak stala jedna z nejuniverzálnějších křivek v dějinách lidstva.  
Parabola spojuje dávné ideály krásy s moderní přesností vědy — a připomíná nám, že matematika není jen nástroj k počítání, ale i způsob, jak se dívat na svět s porozuměním a úžasem.
\end{historicalbox}

\section{Kvadratické funkce}
Zatímco \textbf{lineární funkce} vyjadřuje závislost, při níž se hodnota funkce mění stále stejným způsobem 
(rovnoměrně), \important{kvadratická funkce} nám umožňuje popsat situace, kdy se tato změna zrychluje nebo zpomaluje.  
V praxi se s takovými vztahy setkáváme velmi často — ať už jde o dráhu tělesa při volném pádu, 
tvar paraboly reflektoru nebo zisk firmy, který po určitém bodě začne klesat.

\begin{definitionbox}
\textbf{Kvadratická funkce} je každá funkce, kterou lze vyjádřit předpisem
\[
f(x) = ax^2 + bx + c,
\]
kde \(a, b, c \in \mathbb{R}\) a \(a \neq 0\).

\medskip
Grafem kvadratické funkce je \textbf{plynulá nepřerušovaná křivka}, která se nazývá\textbf{parabola}, 
která je souměrná podle své \textbf{osy souměrnosti} — tedy osa rovnoběžná s osou $y$.  
Definiční obor kvadratické funkce je \(D(f) = \mathbb{R}\), tedy všechny reálné hodnoty proměnné \(x\).
\end{definitionbox}


\begin{notebox}
Kvadratické funkce často používáme i tam, kde nejde o čistě matematický popis: například k 
modelování závislosti ceny na poptávce, množství světla na vzdálenosti od zdroje nebo trajektorie míče při hodu.  
Jejich význam tedy sahá daleko za hranice školních lavic.
\end{notebox}

\newpage
\section{Graf kvadratické funkce}

\begin{definitionbox}
\textbf{Základní kvadratická funkce} má předpis
\[
f(x) = x^2.
\]

Její grafem je křivka zvaná \textbf{parabola}.  
Tento nejjednodušší tvar je výchozím bodem pro všechny ostatní kvadratické funkce, které vznikají posunutím, protažením nebo převrácením této základní paraboly.

\medskip
Parabola je \textbf{sudá} funkce, protože platí \(f(-x) = f(x)\).  
Její graf je tedy souměrný podle osy $y$.  
Zároveň se skládá ze dvou nekonečných \textbf{ramen}, která se symetricky rozbíhají od jednoho zvláštního bodu – \textbf{vrcholu paraboly}.
\end{definitionbox}

\begin{center}
\begin{axisgraph}[
  xmin=-4, xmax=4,
  ymin=-1, ymax=5,
  samples=200,
  width=10cm,
  height=7cm,
  domain=-2.2:2.2,
  title={Graf funkce $y = x^2$}
]
  \addplot[thick, pastelblue!80!black, smooth] {x^2};
  \draw[dashed, pastelorange] (axis cs:0,-1) -- (axis cs:0,5);
  \filldraw[pastelblue!80!black] (axis cs:0,0) circle (2pt) node[above right] {$V[0;0]$};
\end{axisgraph}
\end{center}

\begin{notebox}
Graf kvadratické funkce \(f(x)=x^2\) má několik základních vlastností:
\begin{itemize}
  \item Je \textbf{sudý}, tedy souměrný podle osy $y$.
  \item Je \textbf{klesající} na intervalu \((-\infty, 0)\) a \textbf{rostoucí} na intervalu \((0, +\infty)\).
  \item Není \textbf{prostý} – například \(f(-1) = f(1) = 1\).
  \item Není \textbf{lichý}, protože neplatí \(f(-x) = -f(x)\).
  \item Má \textbf{minimum} ve vrcholu paraboly \(V[0;0]\), kde nabývá nejmenší možné hodnoty \(f(x) = 0\).
  \item \textbf{Nemá maximum}, protože její ramena rostou do nekonečna.
  \item Její \textbf{obor hodnot} je \(H(f) = \langle 0; +\infty)\).
\end{itemize}
\end{notebox}


Každou kvadratickou funkci si můžeme představit jako „přetvořenou“ podobu základní paraboly \(y=x^2\).  
Změníme-li její koeficienty, můžeme ji posunout, zúžit, převrátit nebo přenést jinam do souřadnicové roviny.  
Všechny tyto změny si v následující části ukážeme na obecném tvaru \(f(x) = ax^2 + bx + c\).

\subsection{Různé podoby kvadratických funkcí}

\begin{examplebox}
Podívejme se nyní na několik různých kvadratických funkcí a všimněme si, jak se jejich grafy liší:

\[
\begin{aligned}
a)\quad & y = -x^2\\[4pt]
b)\quad & y = 3x^2\\[4pt]
c)\quad & y = x^2 + 2\\[4pt]
d)\quad & y = x^2 - 4x + 8
\end{aligned}
\]

Každá z těchto funkcí je kvadratická, a tedy její grafem je vždy \textbf{parabola}.  
Na první pohled je ale patrné, že se tyto paraboly chovají různě:
\begin{itemize}
  \item U funkce \(y = -x^2\) je parabola převrácená směrem dolů — hodnota \(a = -1\) tedy způsobila \important{otočení}.
  \item U funkce \(y = 3x^2\) se parabola „zúžila“ — větší koeficient \(a\) znamená prudší \important{růst i pokles}
  \item U funkce \(y = x^2 + 2\) se celý graf posunul \important{směrem nahoru}, ale tvar zůstal stejný.
  \item U funkce \(y = x^2 - 4x + 8\) se parabola posunula nejen vzhůru, ale i směrem doprava – a její vrchol už neleží v počátku.
\end{itemize}
\end{examplebox}

\begin{center}
\begin{multicols}{2}
\begin{axisgraph}[
  xmin=-3, xmax=3,
  ymin=-5, ymax=2,
  samples=200,
  width=9cm, height=7cm,
  title={a) $y = -x^2$}
]
  \addplot[thick, pastelblue!80!black, smooth, domain=-2:2] {-x^2};
\end{axisgraph}

\begin{axisgraph}[
  xmin=-3, xmax=3,
  ymin=-2, ymax=5,
  samples=200,
  width=9cm, height=7cm,
  title={b) $y = 5x^2$}
]
  \addplot[thick, pastelblue!80!black, smooth, domain=-1.5:1.5] {5*x^2};
\end{axisgraph}

\begin{axisgraph}[
  xmin=-3, xmax=3,
  ymin=-1, ymax=6,
  samples=200,
  width=9cm, height=7cm,
  title={c) $y = x^2 + 2$}
]
  \addplot[thick, pastelblue!80!black, smooth, domain=-2:2] {x^2 + 2};
\end{axisgraph}

\begin{axisgraph}[
  xmin=-1, xmax=5,
  ymin=-1, ymax=9,
  samples=200,
  width=9cm, height=7cm,
  title={d) $y = x^2 - 4x + 8$}
]
  \addplot[thick, pastelblue!80!black, smooth, domain=-1:6] {x^2 - 4*x + 8};
\end{axisgraph}
\end{multicols}
\end{center}

Na těchto čtyřech grafech je dobře vidět, jak se mění poloha a tvar paraboly podle hodnot koeficientů \(a\), \(b\) a \(c\).  
Zatímco parametr \(a\) určuje, zda je parabola otevřená nahoru či dolů a jak je „široká“, koeficient \(c\) ji posouvá svisle a \(b\) má vliv na její vodorovné posunutí.  
Změnou těchto tří čísel tedy dokážeme parabolu v rovině zrcadlit, rozšířit, posunout nebo převrátit, aniž bychom změnili její skutečný tvar.

\medskip
Na rozdíl od lineární funkce, jejíž grafem je přímka, má kvadratická funkce tvar \textbf{křivky} s charakteristickým bodem, který nazýváme \textbf{vrchol paraboly}.  
Právě tento bod představuje jakýsi pomyslný „střed“ křivky — místo, kde se mění směr růstu na klesání nebo naopak.

\begin{notebox}
Z hlediska geometrie jsou však \textbf{všechny paraboly shodné}.  
Každou parabolu lze z té základní, $y=x^2$ získat pouhým posunutím, zrcadlením nebo roztažením v rovině, aniž by se změnil její tvar.  
Neexistují tedy „různé druhy“ parabol – liší se pouze polohou, směrem otevření a měřítkem, nikoliv samotným tvarem.
\end{notebox}

\subsection{Vrchol kvadratické funkce a osa souměrnosti}

Už víme, že graf kvadratické funkce má vždy tvar paraboly a že právě její vrchol určuje, kde se křivka „láme“.  
Nyní si ukážeme, jak lze souřadnice tohoto vrcholu \([p; q]\) vypočítat přímo z obecného tvaru funkce.

\medskip
Začněme funkcí:
\[
f(x) = ax^2 + bx + c.
\]
Chceme najít takovou hodnotu \(x_v\), pro kterou má funkce nejmenší (nebo největší) hodnotu.  
Všimněme si, že člen \(ax^2\) vždy roste (nebo klesá) symetricky, zatímco výraz \(bx\) parabolu „naklání“ na jednu stranu.  
Abychom zjistili, kde se parabola obrací, potřebujeme tedy nalézt bod, ve kterém se její růst a pokles vyrovnají.

\medskip
Z geometrického pohledu je to \textbf{osa souměrnosti paraboly} — přímka, která prochází jejím vrcholem.  
Pro libovolnou kvadratickou funkci má rovnici:
\[
x = -\frac{b}{2a}.
\]
Tento vztah nám dává právě souřadnici \(x_v\) vrcholu.

\medskip
Jakmile známe \(x_v\), můžeme snadno spočítat i odpovídající hodnotu \(y_v\) dosazením do původního předpisu:
\[
y= ax^2 + bx +c,
\]
kde za $x$ dosadíme $-\dfrac{b}{2a}$ z čehož nám poté vznikne:
\[
y_v = f(x_v) = a\left(-\frac{b}{2a}\right)^2 + b\left(-\frac{b}{2a}\right) + c.
\]
Po úpravě dostáváme:
\[
y_v = \frac{4ac - b^2}{4a} = \frac{4ac}{4a} - \frac{b^2}{4a} = c - \frac{b^2}{4a}.
\]

\medskip
Tedy vrchol paraboly \(V[x_v; y_v]\) má obecně souřadnice:
\[
x_v = -\frac{b}{2a}, \qquad y_v = c - \frac{b^2}{4a}.
\]

\begin{definitionbox}
\textbf{Vrchol paraboly a osa souměrnosti}

Pro kvadratickou funkci
\[
f(x) = ax^2 + bx + c, \quad a \neq 0,
\]
platí:
\[
x_v = -\frac{b}{2a}, \qquad y_v = c - \frac{b^2}{4a}.
\]

\begin{itemize}
  \item Vrchol paraboly má souřadnice \(V[x_v; y_v]\).
  \item Osa souměrnosti má rovnici \(x = x_v\).
  \item Pokud \(a > 0\), parabola má v bodě \(V\) \textbf{minimum},  
        pokud \(a < 0\), má v bodě \(V\) \textbf{maximum}.
\end{itemize}
\end{definitionbox}

\medskip
Na konkrétních příkladech si nyní ukážeme, jak tyto hodnoty určit v praxi — a poté si vysvětlíme jiný, elegantní způsob, jak k nim dojít, tzv. \textbf{doplněním na čtverec}.  
Tento způsob navíc vede k přehlednému tvaru
\[
f(x) = a(x - p)^2 + q,
\]
ze kterého lze souřadnice vrcholu \([p; q]\) ihned vyčíst, aniž bychom museli počítat pomocí vzorců.
\begin{examplebox}
Určete souřadnice vrcholu paraboly kvadratické funkce
\[
f(x) = x^2 - 4x + 8.
\]

Z obecného tvaru \(f(x) = ax^2 + bx + c\) vyplývá, že:
\[
a = 1, \quad b = -4, \quad c = 8.
\]

Dosadíme do vzorců pro vrchol:
\[
x_v = -\frac{b}{2a} = -\frac{-4}{2 \cdot 1} = 2,
\]
\[
y_v = c - \frac{b^2}{4a} = 8 - \frac{(-4)^2}{4 \cdot 1} = 8 - 4 = 4.
\]

Vrchol paraboly tedy má souřadnice:
\[
V[2; 4],
\]
a osa souměrnosti má rovnici \(x = 2\).
\end{examplebox}

\begin{center}
\begin{axisgraph}[
  xmin=-1, xmax=6,
  ymin=-1, ymax=10,
  samples=200,
  width=9cm,
  height=6cm,
  title={Graf funkce $y = x^2 - 4x + 8$}
]
  \addplot[thick, pastelblue!80!black, smooth, domain=-1:6] {x^2 - 4*x + 8};
  \draw[dashed, gray] (axis cs:2,-1) -- (axis cs:2,10);
  \filldraw[pastelblue!80!black] (axis cs:2,4) circle (2pt) node[below right] {$V[2;4]$};
\end{axisgraph}
\end{center}

Jak vidíme na grafu, parabola se skutečně otevírá směrem nahoru (\(a>0\))  
a její vrchol se nachází v bodě \(V[2;4]\), tedy dvě jednotky vpravo a čtyři jednotky nahoru od počátku.

\medskip
Pomocí vzorců jsme tedy dokázali přesně určit, kde se parabola „láme“.  
V praxi však můžeme těchto hodnot dosáhnout i bez zapamatování vzorců — stačí vhodně upravit samotný předpis funkce.  
Této úpravě se říká \important{doplnění na čtverec}.

\subsection{Vrcholový tvar kvadratické funkce – doplnění na čtverec}

Doplnění na čtverec nám umožní přepsat kvadratickou funkci do přehlednějšího tvaru:
\[
f(x) = a(x - p)^2 + q,
\]
kde \([p; q]\) jsou souřadnice vrcholu paraboly.  
Postup si ukážeme na našem příkladu \(f(x) = x^2 - 4x + 8\).

\medskip
Nejprve se zaměříme na část s proměnnou \(x\):
\[
x^2 - 4x.
\]
Abychom z ní vytvořili dokonalý čtverec, připomeneme si vztah:
\[
(a + b)^2 = a^2 + 2ab + b^2, \quad (a - b)^2 = a^2 - 2ab + b^2,
\]
kterého využijeme zde. 
\[
x^2 - 4x + 4 = (x - 2)^2 
\]
Vidíme, že náš výraz však neodpovídá původnímu zadání:
\[
x^2 - 4x + 4 \neq x^2 - 4x + 8
\]
Oproti původnímu zadání tam máme člen $c$ o $4$ větší. Takže ho taky musíme přičíst:

\[
(x^2 - 4x + 4) + 4 = x^2 - 4x + 8
\]
A protože nyní již známe hodnotu prvního členu v závorce $x^2 - 4x + 4 = (x - 2)^2 $, můžeme jej tedy nahradit a dostaneme:
\[
(x - 2)^2 + 4 = x^2 - 4x + 8
\]
Čehož jsme chtěli dosáhnou a z tohoto tvaru, jež se nazývá \important{Vrcholový tvar} můžeme tedy vyčíst hodnoty vrcholu:
\[
V=[2;4]
\]

\medskip
Tento způsob je nejen elegantní, ale také názorný — studenti často ihned pochopí, co znamená posun \(p\) a \(q\) v předpisu, a mohou graf snadno odhadnout bez počítání.

\begin{examplebox}
\textbf{Příklad 1:}  
Určete vrcholový tvar funkce \(f(x) = 2x^2 - 8x + 5.\)

\medskip
Nejprve si z prvních dvou členů vytkneme číslo \(2\), aby zbyl výraz se součinitelem u \(x^2\) roven jedné:
\[
f(x) = 2(x^2 - 4x) + 5.
\]
Nyní doplníme na čtverec:
\[
x^2 - 4x + 4 = (x - 2)^2.
\]
Protože však my nahrazujeme $x^2 - 4$ a opět tam máme navíc parametr $c$ navíc, tak ho taky musíme odečíst:
\[
(x^2 - 4x + 4) - 4 = (x - 2)^2 - 4.
\]
\[
x^2 - 4x = (x - 2)^2 - 4.
\]
A nyní můžeme dosadit z původní upravené funkce $f(x) = 2(x^2 - 4x) + 5$ za výraz $(x^2 - 4x)$ náš čtverec $(x - 2)^2 - 4$.
Tudíž poté máme funkci:
\[
f(x) = 2\big[(x - 2)^2 - 4\big] + 5.
\]
Roznásobíme hranatou závorku a upravíme:
\[
f(x) = 2(x - 2)^2 - 8 + 5 = 2(x - 2)^2 - 3.
\]
Vrcholový tvar funkce tedy je:
\[
f(x) = 2(x - 2)^2 - 3,
\]
a vrchol má souřadnice \(V[2; -3].\)

\medskip
Vidíme, že větší hodnota \(a = 2\) parabolu „zúžila“, ale tvar zůstal stejný.  
Parabola je otevřená směrem nahoru a má minimum v bodě \(V[2; -3].\)
\end{examplebox}

\begin{center}
\begin{axisgraph}[
  xmin=-1, xmax=5,
  ymin=-4, ymax=5,
  samples=200,
  width=9cm,
  height=7cm,
  title={$y = 2x^2 - 8x + 5 = 2(x - 2)^2 - 3$}
]
  \addplot[thick, pastelblue!80!black, smooth, domain=-1:5] {2*(x - 2)^2 - 3};
  \draw[dashed, gray] (axis cs:2,-6) -- (axis cs:2,5);
  \filldraw[pastelblue!80!black] (axis cs:2,-3) circle (2pt) node[below] {$V[2;-3]$};
\end{axisgraph}
\end{center}

\begin{examplebox}
\textbf{Příklad 2:}  
Určete vrcholový tvar funkce \(f(x) = -x^2 + 6x - 5.\)

\medskip
Nejprve vytkneme záporné znaménko, aby před \(x^2\) stála jednička:
\[
f(x) = -\big(x^2 - 6x + 5\big).
\]
Nyní se zaměříme na část v závorce:
\[
x^2 - 6x = (x - 3)^2 - 9.
\]
Dosadíme zpět:
\[
f(x) = -\big[(x - 3)^2 - 9 + 5\big] = -\big[(x - 3)^2 - 4\big].
\]
Roznásobíme mínus:
\[
f(x) = -(x - 3)^2 + 4.
\]
Vrcholový tvar tedy je:
\[
f(x) = -(x - 3)^2 + 4,
\]
a vrchol má souřadnice \(V[3; 4].\)

\medskip
Protože \(a = -1\), parabola je otevřená směrem dolů a v bodě \(V[3; 4]\) má \textbf{maximum.}
\end{examplebox}

\begin{center}
\begin{axisgraph}[
  xmin=-1, xmax=6,
  ymin=-2, ymax=6,
  samples=200,
  width=9cm,
  height=6cm,
  title={$y = -x^2 + 6x - 5 = -(x - 3)^2 + 4$}
]
  \addplot[thick, pastelblue!80!black, smooth, domain=-1:7] {-(x - 3)^2 + 4};
  \draw[dashed, gray] (axis cs:3,-3) -- (axis cs:3,6);
  \filldraw[pastelblue!80!black] (axis cs:3,4) circle (2pt) node[below right] {$V[3;4]$};
\end{axisgraph}
\end{center}

\medskip
Z obou příkladů je patrné, že doplnění na čtverec je univerzální metoda – funguje bez ohledu na znaménko i velikost \(a\).  
Získaný vrcholový tvar \(f(x) = a(x - p)^2 + q\) má nejen jasný geometrický význam, ale zároveň umožňuje snadno číst polohu a vlastnosti paraboly přímo z rovnice.


\subsection{Průsečíky paraboly s osami}

Kromě vrcholu nás u grafu kvadratické funkce často zajímají také její \textbf{průsečíky s osami}.  
Pomáhají nám při náčrtu paraboly, protože určují, kde graf protíná osy souřadnicové soustavy.

\subsubsection*{Průsečík s osou \(y\)}

Průsečík s osou \(y\) je bod, kde má parabola souřadnici \(x = 0\).  
Stačí tedy do předpisu funkce dosadit \(x = 0\) a vypočítat odpovídající hodnotu \(y\):
\[
y = a \cdot 0^2 + b \cdot 0 + c = c.
\]
Z toho plyne, že parabola \textbf{vždy protíná osu \(y\) v bodě \([0; c]\)}.  
Tento průsečík existuje pro každou kvadratickou funkci, protože osa \(y\) leží v jejím definičním oboru.

\begin{examplebox}
\textbf{Příklad:}  
Určete průsečík paraboly \(f(x) = 2x^2 - 4x + 3\) s osou \(y\).

Dosadíme \(x = 0\):
\[
f(0) = 2 \cdot 0^2 - 4 \cdot 0 + 3 = 3.
\]
Průsečík je tedy bod \([0; 3]\).
\end{examplebox}

\subsubsection*{Průsečíky s osou \(x\)}

Průsečíky s osou \(x\) odpovídají místům, kde má funkce hodnotu \(y = 0\).  
To znamená, že hledáme řešení rovnice
\[
ax^2 + bx + c = 0.
\]
Tato rovnice se nazývá \textbf{kvadratická rovnice} a budeme se jí věnovat v samostatné kapitole.  
Zatím si jen uvedeme, že její řešení (pokud existují) lze vyjádřit pomocí vzorce:
\[
x_{1,2} = \frac{-b \pm \sqrt{b^2 - 4ac}}{2a}.
\]
Hodnoty \(x_1\) a \(x_2\) představují souřadnice průsečíků paraboly s osou \(x\).

\medskip
Podle hodnoty výrazu pod odmocninou — tzv. \textbf{diskriminantu}
\[
D = b^2 - 4ac,
\]
může mít parabola:
\begin{itemize}
  \item dva průsečíky s osou \(x\) (\(D > 0\)),
  \item jeden průsečík — dotýká se osy \(x\) (\(D = 0\)),
  \item nebo žádný průsečík (\(D < 0\)).
\end{itemize}

\begin{definitionbox}
\textbf{Shrnutí: průsečíky paraboly s osami}
\begin{itemize}
  \item \textbf{Průsečík s osou \(y\)}: \([0; c]\).
  \item \textbf{Průsečíky s osou \(x\)}: \([x_1; 0]\), \([x_2; 0]\),  
  kde \(x_{1,2} = \dfrac{-b \pm \sqrt{b^2 - 4ac}}{2a}\).
  \item Počet průsečíků s osou \(x\) závisí na hodnotě diskriminantu \(D = b^2 - 4ac\).
\end{itemize}
\end{definitionbox}

\begin{center}
\begin{axisgraph}[
  xmin=-3, xmax=5,
  ymin=-5, ymax=5,
  samples=200,
  width=10cm,
  height=10cm,
  title={Graf funkce $y = x^2 - 2x - 3$}
]
  \addplot[thick, pastelblue!80!black, smooth, domain=-4:5] {x^2 - 2*x - 3};
  \filldraw[pastelblue!80!black] (axis cs:-1,0) circle (2pt) node[above] {$A[-1;0]$};
  \filldraw[pastelblue!80!black] (axis cs:3,0) circle (2pt) node[above] {$B[3;0]$};
  \filldraw[pastelblue!80!black] (axis cs:0,-3) circle (2pt) node[below left] {$C[0;-3]$};

\end{axisgraph}
\end{center}

Z grafu vidíme, že parabola protíná osu \(y\) v bodě \(C[0; -3]\)  
a osu \(x\) ve dvou bodech \(A[-1; 0]\) a \(B[3; 0]\).  
Tento příklad odpovídá situaci, kdy má diskriminant \(D > 0\) — tedy dvě různá reálná řešení.

\begin{examplebox}
\textbf{Příklad:}  
Podívejme se, jak se mění tvar a poloha paraboly podle hodnoty diskriminantu \(D = b^2 - 4ac\).

\begin{multicols}{3}
\begin{axisgraph}[
  xmin=-3, xmax=3,
  ymin=-2, ymax=5,
  samples=200,
  width=7cm, height=6cm,
  title={$D > 0$ – dva průsečíky}
]
  % y = x^2 - 1 → D = 4
  \addplot[thick, pastelblue!80!black, smooth, domain=-3:3] {x^2 - 1};
  \filldraw[pastelblue!80!black] (axis cs:-1,0) circle (1.5pt);
  \filldraw[pastelblue!80!black] (axis cs:1,0) circle (1.5pt);
  \node[below] at (axis cs:-1,0) {\small $x_1$};
  \node[below] at (axis cs:1,0) {\small $x_2$};
\end{axisgraph}

\begin{axisgraph}[
  xmin=-2, xmax=4,
  ymin=-2, ymax=5,
  samples=200,
  width=7cm, height=6cm,
  title={$D = 0$ – jeden průsečík}
]
  % y = (x-1)^2 → D = 0
  \addplot[thick, pastelblue!80!black, smooth, domain=-2:4] {(x-1)^2};
  \filldraw[pastelblue!80!black] (axis cs:1,0) circle (1.5pt);
  \node[above] at (axis cs:1,0) {\small $x_1=x_2$};
\end{axisgraph}

\begin{axisgraph}[
  xmin=-3, xmax=3,
  ymin=-2, ymax=5,
  samples=200,
  width=7cm, height=6cm,
  title={$D < 0$ – žádný průsečík}
]
  % y = x^2 + 3 → D < 0
  \addplot[thick, pastelblue!80!black, smooth, domain=-3:3] {x^2 + 3};
\end{axisgraph}
\end{multicols}

\medskip
Z těchto tří grafů vidíme, že:
\begin{itemize}
  \item Pokud \(D > 0\), parabola protíná osu \(x\) ve dvou bodech – má dva reálné kořeny.
  \item Pokud \(D = 0\), parabola se osy \(x\) dotýká – má právě jeden kořen (dvojnásobný).
  \item Pokud \(D < 0\), parabola osu \(x\) neprotíná – nemá žádný reálný kořen.
\end{itemize}
\end{examplebox}


\subsection{Souvislost s analytickou geometrií}

Kvadratické funkce tvoří přirozený most mezi algebrou a geometrií.  
V této kapitole jsme se zabývali pouze funkcí jedné proměnné, jejíž grafem je parabola otevřená buď nahoru, nebo dolů.  
V pozdější části nazvané \textbf{Analytická geometrie} se k parabole znovu vrátíme – tentokrát už jako k jedné z tzv. \textbf{kuželoseček}.

\medskip
Kuželosečky představují skupinu křivek, které vznikají průnikem roviny s kuželem.  
Patří mezi ně:
\[
\text{elipsa, parabola a hyperbola.}
\]
Zjistíme, že parabola je pouze jeden speciální případ těchto křivek a že její rovnice
\[
y = ax^2 + bx + c
\]
je jen jednodušším zápisem obecnějšího kvadratického vztahu mezi proměnnými \(x\) a \(y\).  
Z pohledu geometrie se tedy kvadratická funkce stane prvním krokem k pochopení celé rodiny kuželoseček.

\begin{notebox}
V analytické geometrii se naučíme nejen parabolu popsat pomocí rovnic, ale také určit její ohnisko, přímku řídicí a odvodit z nich všechny její základní vlastnosti.  
To, co jsme si zde ukázali pro kvadratickou funkci, tak tvoří základ, na který budeme později navazovat.
\end{notebox}

\subsection{Parabola v praxi a kolem nás}

Paraboly nejsou jen součástí učebnic, ale i našeho každodenního života.  
Jejich charakteristická vlastnost – soustředit paprsky nebo zvuk do jednoho bodu – se využívá v celé řadě technických i architektonických aplikací.

\begin{historicalbox}
\textbf{Příklady praktického využití paraboly:}

\begin{itemize}
  \item \textbf{Reflektory a světlomety} – parabola odráží světelné paprsky do rovnoběžného svazku, což umožňuje rovnoměrné osvětlení vozovek.
  \item \textbf{Satelitní antény a teleskopy} – parabolické zrcadlo soustřeďuje signály do ohniska, kde je zachytí přijímač.
  \item \textbf{Akustické paraboly} – směrují a zesilují zvuk z jednoho směru (např. na sportovních stadionech nebo u mikrofonů).
  \item \textbf{Architektura a stavitelství} – oblouky, mosty či střešní konstrukce využívají parabolický tvar pro rozložení sil a estetický efekt.
  \item \textbf{Fyzika a mechanika} – dráha tělesa vrženého do vzduchu je parabolická, stejně jako proud vody z fontány.
\end{itemize}

\medskip
Parabola tak propojuje svět matematiky, fyziky i techniky.  
Z jednoduché rovnice \(y = ax^2 + bx + c\) se stává univerzální model přírodních i člověkem vytvořených jevů.
\end{historicalbox}



\subsection{Určení předpisu kvadratické funkce ze tří bodů}

Kvadratická funkce má obecný tvar
\[
f(x) = ax^2 + bx + c.
\]
Z jejích tří koeficientů \(a\), \(b\), \(c\) lze jednoznačně určit graf – a opačně:  
známe-li tři body, kterými parabola prochází, můžeme z nich určit právě tyto tři neznámé.


\subsubsection*{Postup}

Máme-li zadány tři body:
\[
A[x_1; y_1], \quad B[x_2; y_2], \quad C[x_3; y_3],
\]
dosadíme jejich souřadnice do obecného tvaru \(f(x) = ax^2 + bx + c\).  
Získáme tak soustavu tří lineárních rovnic o třech neznámých:
\[
\begin{cases}
a x_1^2 + b x_1 + c = y_1,\\
a x_2^2 + b x_2 + c = y_2,\\
a x_3^2 + b x_3 + c = y_3.
\end{cases}
\]
Tuto soustavu následně vyřešíme (např. dosazením, sčítací metodou nebo maticově) a nalezené hodnoty \(a\), \(b\), \(c\) dosadíme zpět do předpisu funkce.


\begin{examplebox}
\textbf{Příklad:}  
Najděte předpis kvadratické funkce, která prochází body
\[
A[0; 2], \quad B[1; 1], \quad C[2; 6].
\]

\textbf{Řešení:}

Dosadíme do obecného tvaru \(y = ax^2 + bx + c\):
\[
\begin{cases}
a \cdot 0^2 + b \cdot 0 + c = 2 \\[2pt]
a \cdot 1^2 + b \cdot 1 + c = 1 \\[2pt]
a \cdot 2^2 + b \cdot 2 + c = 6
\end{cases}
\quad \Rightarrow \quad
\begin{cases}
c = 2 \\[2pt]
a + b + 2 = 1 \\[2pt]
4a + 2b + 2 = 6
\end{cases}
\]

Z druhé rovnice: \(a + b = -1\).  
Z třetí: \(4a + 2b = 4 \Rightarrow 2a + b = 2.\)

Odečteme obě rovnice:
\[
(2a + b) - (a + b) = 2 - (-1) \quad \Rightarrow \quad a = 3.
\]
Dosadíme zpět: \(3 + b = -1 \Rightarrow b = -4.\)

Z první rovnice víme \(c = 2.\)

\[
\boxed{f(x) = 3x^2 - 4x + 2.}
\]

\begin{center}
\begin{axisgraph}[
  xmin=-1, xmax=3,
  ymin=-1, ymax=7,
  samples=200,
  width=9cm,
  height=6cm,
  title={Parabola procházející body $A[0;2], B[1;1], C[2;6]$}
]
  \addplot[thick, pastelblue!80!black, smooth, domain=-1:3] {3*x^2 - 4*x + 2};
  \addplot[only marks, mark=*] coordinates {(0,2) (1,1) (2,6)};
  \node[below right] at (axis cs:0,2) {$A$};
  \node[below] at (axis cs:1,1) {$B$};
  \node[above left] at (axis cs:2,6) {$C$};
\end{axisgraph}
\end{center}
\end{examplebox}


\begin{notebox}
Tento postup se používá nejen v matematice, ale i v praxi — např. při určování trajektorie pohybu, modelování oblouků, interpolaci dat nebo při grafickém návrhu objektů, které mají parabolický tvar (např. mostní oblouky, reflektory, satelitní antény).
\end{notebox}

\begin{notebox}
Pokud je zadáno více než tři body, nelze obecně najít jedinou přesnou kvadratickou funkci, která by procházela všemi — používá se tzv. \textbf{aproximace} nebo \textbf{regresní parabola}, což je téma, ke kterému se vrátíme později v části o zpracování dat.
\end{notebox}






\section{Kavdratické funkce v nerovnosti}

Každou kvadratickou funkci jsme dosud znázorňovali jako křivku — tedy množinu bodů, které splňují rovnici
\[
y = f(x) = ax^2 + bx + c.
\]
Pokud však místo rovnosti použijeme nerovnost, například
\[
y > x^2 \quad \text{nebo} \quad y \le x^2 + 3x - 2,
\]
už nepopisujeme jen křivku, ale \textbf{celou oblast} v rovině.

\medskip
Vztah \(y > f(x)\) znamená, že hledáme všechny body, jejichž souřadnice \(y\) je \textbf{větší} než hodnota funkce v daném bodě \(x\).  
Jinými slovy, jde o všechny body \textbf{nad grafem paraboly}.  
Naopak, pro \(y < f(x)\) leží body \textbf{pod grafem}.  
Pokud použijeme znaménka \(\ge\) nebo \(\le\), k oblasti zároveň patří i body samotné na parabole.

\begin{examplebox}
\textbf{Příklad:}  
Rozhodněte, jakou oblast v rovině popisuje nerovnost
\[
y \le x^2.
\]

Tato nerovnost zahrnuje všechny body, jejichž souřadnice \(y\) je menší nebo rovna hodnotě \(x^2\).  
Protože graf funkce \(y = x^2\) je parabola otevřená nahoru, bude hledaná oblast \textbf{uvnitř paraboly}, tedy „pod jejím obloukem“.  
Bod \((0,0)\) na parabole sem patří (protože \(\le\) znamená i rovnost).

\begin{center}
\begin{axisgraph}[
  xmin=-3, xmax=3,
  ymin=-1, ymax=5,
  samples=200,
  width=9cm,
  height=6cm,
  title={Oblast $y \le x^2$}
]
  \addplot[domain=-3:3, thick, pastelblue!80!black, smooth] {x^2};
  \addplot[domain=-3:3, draw=none, fill=pastelblue!50, opacity=0.75] {x^2} \closedcycle;
\end{axisgraph}
\end{center}
\end{examplebox}


\medskip
\noindent
Obecně tedy platí:

\begin{itemize}
  \item \(y > f(x)\): oblast \textbf{nad grafem} (včetně „vnitřku“ pro parabolu otevřenou dolů),
  \item \(y < f(x)\): oblast \textbf{pod grafem} (včetně „vnitřku“ pro parabolu otevřenou nahoru),
  \item \(y \ge f(x)\) nebo \(\le f(x)\): oblast včetně hranice – paraboly samotné.
\end{itemize}

\begin{notebox}
Tyto nerovnosti mezi \(y\) a \(f(x)\) se často používají při popisu geometrických tvarů — například při určování stínů, trajektorií, polí nebo oblastí v analytické geometrii.  
Ukazují, že kvadratická funkce neurčuje jen křivku, ale může vymezovat i prostor kolem ní.
\end{notebox}


\subsection{Kvadratické funkce s absolutní hodnotou}

Absolutní hodnota v předpisu kvadratické funkce může výrazně změnit její tvar.  
Pojďme si ukázat tři základní případy, se kterými se budeme setkávat.


\subsubsection*{1. Typ \(\displaystyle y = \lvert ax^2 + bx + c\rvert\)}

V tomto případě je absolutní hodnota aplikována až na výsledek celé funkce, tedy na hodnotu \(y\).  
Platí, že:
\[
|y| =
\begin{cases}
y, & \text{pro } y \ge 0,\\
-y, & \text{pro } y < 0.
\end{cases}
\]

Proto postupujeme takto:
1. Nejprve sestrojíme graf funkce bez absolutní hodnoty – \(y = ax^2 + bx + c\).
2. Poté všechny části, které leží pod osou \(x\), „překlopíme“ osově souměrně podle osy \(x\).

\begin{examplebox}
\textbf{Příklad:}  
Sestrojte graf funkce \(y = |x^2 - 4|\).

\begin{center}
\begin{axisgraph}[
  xmin=-3, xmax=3,
  ymin=-1, ymax=5,
  samples=200,
  width=9cm,
  height=6cm,
  title={$y = |x^2 - 4|$}
]
  % původní parabola
  \addplot[dashed, gray!60, domain=-3:3, smooth] {x^2 - 4};
  % překlopená část pod osou x
  \addplot[thick, pastelblue!80!black, domain=-3:3, smooth]
    {abs(x^2 - 4)};
  % pomocná osa
  \draw[gray, thin] (axis cs:-3,0)--(axis cs:3,0);
\end{axisgraph}
\end{center}

Části původní paraboly pod osou \(x\) se „odrazí“ nahoru – vznikne tvar připomínající dvě spojené paraboly.
\end{examplebox}


\subsubsection*{2. Typ \(\displaystyle y = ax^2 + b|x| + c\)}

Absolutní hodnota se zde nachází pouze u proměnné \(x\).  
Tím se celý graf stává \textbf{sudou funkcí}, a tedy \textbf{souměrným podle osy \(y\)}.

Postup:
1. Nejprve sestrojíme graf funkce bez absolutní hodnoty (např. jen pro \(x \ge 0\)).
2. Poté tuto část osově zobrazíme podle osy \(y\).

\begin{examplebox}
\textbf{Příklad:}  
Sestrojte graf funkce \(y = x^2 + 2|x| - 3.\)

\begin{center}
\begin{axisgraph}[
  xmin=-4, xmax=4,
  ymin=-4, ymax=5,
  samples=200,
  width=9cm,
  height=6cm,
  title={$y = x^2 - 2|x| - 3$}
]
  \addplot[thick, pastelblue!80!black, domain=-4:4, smooth]
    {x^2 - 2*abs(x) - 3};
  \draw[dashed, gray!60, thin] (axis cs:0,-4)--(axis cs:0,5);
\end{axisgraph}
\end{center}

Graf je zřetelně souměrný podle osy \(y\) a má tvar „rozšířené“ paraboly s dolíkem uprostřed.
\end{examplebox}


\subsubsection*{3. Typ \(\displaystyle y = \lvert ax^2 + b|x| + c\rvert\)}

Tento typ kombinuje oba předchozí případy.  
Nejdříve se absolutní hodnota promítne do proměnné \(x\) (vzniká sudá funkce), a poté se absolutní hodnota uplatní i na výsledek.  
Tedy nejprve se parabola stane souměrnou podle osy \(y\), a pak se všechny části pod osou \(x\) zrcadlí nahoru.

\begin{examplebox}
\textbf{Příklad:}  
Sestrojte graf funkce \(y = \lvert -2x^2 + 4|x| + 1\rvert.\)

\begin{center}
\begin{axisgraph}[
  xmin=-3, xmax=3,
  ymin=-1, ymax=7,
  samples=200,
  width=9cm,
  height=6cm,
  title={$y = |-2x^2 + 4|x| + 1|$}
]
  \addplot[thick, pastelblue!80!black, domain=-3:3, smooth]
    {abs(-2*x^2 + 4*abs(x) + 1)};
  \draw[dashed, gray!60, thin] (axis cs:0,-1)--(axis cs:0,7);
  \draw[gray!60, thin] (axis cs:-3,0)--(axis cs:3,0);
\end{axisgraph}
\end{center}

Tento typ má často tvar připomínající „dvě spojené kopule“ — funkce je zároveň sudá a nezáporná.
\end{examplebox}



\begin{notebox}
Kvadratické funkce s absolutní hodnotou jsou krásným příkladem kombinace symetrie a nelinearity.  
V analytické geometrii se s nimi znovu setkáme při popisu částí kuželoseček a při konstrukci obalových křivek.
\end{notebox}




% -----------------------------------------------------------------
% ---------------------- P Ř Í K L A D Y --------------------------
% -----------------------------------------------------------------

\exercisepart
\setcounter{section}{1}

% =========================
% ÚLOHA 1
% =========================
\begin{exercise}[Tabulka hodnot základní paraboly]
Uvažujme kvadratickou funkci \(f: y=x^2\).  
Zapište do tabulky funkční hodnoty v bodech \(-3;\,-2;\,-1;\,-0{,}5;\,0;\,0{,}5;\,1;\,2;\,3\).

\medskip
\begin{center}
\renewcommand{\arraystretch}{1,25}
\begin{tabular}{|c|*{9}{c|}}
\hline
$x$ & $-3$ & $-2$ & $-1$ & $-0{,}5$ & $0$ & $0{,}5$ & $1$ & $2$ & $3$ \\
\hline
$f(x)=x^2$ &  &  &  &  &  &  &  &  &  \\
\hline
\end{tabular}
\end{center}
\end{exercise}

% =========================
% ÚLOHA 2
% =========================
\begin{exercise}[Srovnání šířky paraboly v jedné soustavě]
Načrtněte graf funkce \(f: y=x^2\) a v téže soustavě souřadnic zakreslete graf funkce
\(f_1: y=\tfrac{1}{3}x^2\). Krátce okomentujte rozdíl.
\end{exercise}

% =========================
% ÚLOHA 3
% =========================
\begin{exercise}[Vrcholový tvar – přímé rýsování]
Sestrojte graf funkce \(g: y=\tfrac{3}{4}(x+1)^2-3\).  
Uveďte vrchol a rovnici osy souměrnosti.
\end{exercise}

% =========================
% ÚLOHA 4
% =========================
\begin{exercise}[Obecný tvar – převod / náčrt]
Sestrojte graf funkce \(h: y=\tfrac{3}{4}x^2+\tfrac{3}{2}x-\tfrac{9}{4}\).  
Doporučení: převeďte do vrcholového tvaru a poté rýsujte.
\end{exercise}

% =========================
% ÚLOHA 5
% =========================
\begin{exercise}[Rodiny posunů v jedné soustavě]
Načrtněte v téže soustavě souřadnic následující grafy:

\begin{tasks}(1)
  \task \(y=x^2+c,\quad c\in\{0;\,3;\,2;\,-0{,}5\}\)
  \task \(y=(x-k)^2,\quad k\in\{0;\,1;\,2{,}5;\,-2;\,-1{,}5\}\)
  \task \(y=(x-k)^2+m,\) pro dvojice \((k,m)\in\{(1,2),\,(1,-2),\,(1,1{,}5)\}\)
  \task \(y=(x-k)^2+m,\) pro dvojice \((k,m)\in\{(-2,2),\,(1,2),\,(0{,}5,2)\}\)
\end{tasks}
\end{exercise}

% =========================
% ÚLOHA 6
% =========================
\begin{exercise}[Graf kvadratické funkce]
Načrtněte grafy funkcí (určete souřadnice vrcholu, souřadnice průsečíků grafu s osou $x$ a s osou $y$; následně z grafu určete obor funkčních hodnot):

\begin{tasks}(3)
  \task $f_1:\; y = x^2 + 4x + 3$
  \task $f_2:\; y = x^2 - 2x + 3$
  \task $f_3:\; y = x^2 + 1$
  \task $f_4:\; y = -x^2 + 2x + 3$
  \task $f_5:\; y = -x^2 + 6x + 9$
  \task $f_6:\; y = 2x^2 - 3x + 6$
  \task $f_7:\; y = 3x^2 + 6x + 3$
  \task $f_8:\; y = 2x^2 + 4x + 1$
  \task $f_9:\; y = \tfrac{3}{4}x^2 - 3x - \tfrac{7}{2}$
\end{tasks}
\end{exercise}

\newpage
% =========================
% ÚLOHA 7
% =========================
\begin{exercise}[Graf kvadratické funkce s absolutními hodnotami]
Načrtněte grafy funkcí (u každé určete vrchol/vrcholy po intervalech, průsečíky s osou $y$ a případně s osou $x$).

\begin{multicols}{3}

\begin{tasks}
  \task $g_1:\; y = x^2 + 2x - 3$
  \task $g_2:\; y = \lvert x^2 + 2x - 3\rvert$
  \task $g_3:\; y = x^2 + 2\lvert x\rvert - 3$
  \task $g_4:\; y = \big\lvert x^2 + 2\lvert x\rvert - 3\big\rvert$
\end{tasks}

\columnbreak

\begin{tasks}
  \task $h_1:\; y = x^2 - 3x$
  \task $h_2:\; y = x^2 - 3\lvert x\rvert$
  \task $h_3:\; y = x\lvert x - 3\rvert$
  \task $h_4:\; y = \lvert x\rvert (x - 3)$
\end{tasks}

\columnbreak

\begin{tasks}
  \task \mbox{$k_1:\; y = \lvert 2x + 1\rvert - x\lvert x - 4\rvert$}
  \task \mbox{$k_2:\; y = x\lvert x - 2\rvert + \lvert x^2 - 2x\rvert$}
  \task \mbox{$k_3:\; y = \lvert x^2 - x\rvert + \lvert x\rvert - 2$}
  \task \mbox{$k_4:\; y = \lvert 9 - x^2\rvert + \lvert 4 - x^2\rvert - 5$}
\end{tasks}

\end{multicols}
\end{exercise}

% =========================
% ÚLOHA 8
% =========================
\begin{exercise}[Kvadratické funkce s absolutní hodnotou]
Načrtněte grafy těchto funkcí:
\begin{tasks}(2)
  \task $y = x|x|$
  \task $y = 2x + |1 - x^2|$
  \task $y = -2x|x - 3|$
  \task $y = x|5 - x| + 2|x| - |2 - x^2|$
\end{tasks}
\end{exercise}


% =========================
% ÚLOHA 9
% =========================
\begin{exercise}[Náčrt z obecného tvaru]
Načrtněte grafy těchto funkcí (vrchol, osa, průsečíky s osami):
\begin{tasks}(2)
  \task \(y=x^2-2x+3\)
  \task \(y=-x^2-6x-8\)
  \task \(y=2x^2+5x-1\)
  \task \(y=-0{,}5x^2+x+2\)
\end{tasks}
\end{exercise}

% =========================
% ÚLOHA 10 (BONUS)
% =========================
\begin{exercise}[BONUS: Pozitivita pomocí doplnění na čtverec]
Dokažte, že hodnota funkce \(y=x^2-4x+5\) je v každém bodě kladná.
\end{exercise}

% =========================
% ÚLOHA 11
% =========================
\begin{exercise}[Sestavení předpisu z vlastností I]
Funkční předpis kvadratické funkce \(f\) zapište rovnicí, víte-li, že platí:  
\(f\) je sudá na \(\mathbb{R}\), hodnota minima je \(-8\) a jeden z průsečíků grafu s osou \(x\) má souřadnice \([2;0]\).
\end{exercise}

% =========================
% ÚLOHA 12
% =========================
\begin{exercise}[Sestavení předpisu z vlastností II]
Funkční předpis kvadratické funkce \(g\) zapište rovnicí, víte-li, že platí:  
funkce \(g\) pro \(x=2\) nabývá \textbf{maxima}, přičemž hodnota maxima je \(4\), a osu \(y\) protíná graf funkce v bodě \([0;1]\).
\end{exercise}

% =========================
% ÚLOHA 13
% =========================
\begin{exercise}[Sestavení předpisu z vlastností III]
Funkční předpis kvadratické funkce \(h\) zapište rovnicí, víte-li, že platí:  
funkce \(h\) je v intervalu \((-\infty;3)\) \textbf{rostoucí}, v intervalu \(\langle 3;\infty)\) \textbf{klesající}, graf prochází počátkem soustavy souřadnic a hodnota maxima je \(18\).
\end{exercise}

% =========================
% ÚLOHA 14
% =========================
\begin{exercise}[Určení koeficientů z funkční rovnice]
Určete reálná čísla \(a,b\) tak, aby pro funkci \(f(x)=ax^2+bx+5\) platilo:
\[
\forall x\in\mathbb{R}:\quad f(x+1)-f(x)=8x+3.
\]
\end{exercise}

% =========================
% ÚLOHA 13
% =========================
\begin{exercise}[Grafy složených racionálních funkcí]
Načrtněte grafy daných funkcí (nezapomeňte určit definiční obor a asymptoty, jsou-li zřejmé):

\begin{tasks}(2)
  \task $y = \dfrac{\frac{x - 5}{x + 1} + 1}{1 - \frac{2x - 1}{x + 1}}$
  
  \task $y = \dfrac{\dfrac{x^3 - 8}{\frac{x^2 + 4}{x + 2} + \frac{2x}{x + 2}}}{\frac{x^3 + 8}{(x - 2)^2 + 2x}}$
  
  \task $y = \left(x + 2 + \dfrac{2}{x - 1}\right) \left(x - 2 + \dfrac{2}{x + 1}\right)$
  
  \task $y = \dfrac{x^3 - 4x + x^2 - 4}{x^2 - x - 2} \cdot x$
  
  \task $y = \dfrac{12x^3 - 5x^2 - 6x - 1}{4x + 1}$
  
  \task $y = 2x^2 - \dfrac{1}{1 - \frac{x}{1 + x}}$
\end{tasks}
\end{exercise}


% =========================
% ÚLOHA 14
% =========================
\begin{exercise}[Určení kvadratické funkce ze tří bodů]
Najděte předpis kvadratické funkce, která prochází třemi zadanými body.

\begin{tasks}(2)
  \task $A[0;1],\; B[-2;5],\; C[2;5]$
  \task $D[4;9],\; E[1;0],\; F[-3;16]$
  \task $G[0;2],\; H[-1;0],\; I[0{,}5;4{,}5]$
  \task $J[-1;5],\; K[0;7],\; L[2;17]$
  \task $M[2;25],\; N[-3;25],\; O[0;1]$
  \task $P[-2;18{,}4],\; Q[2;54{,}4],\; R[0;-3{,}6]$
  \task $S[0;0],\; T[3;45],\; U[-1;-7]$
  \task $V[2;-48],\; W[4;0],\; X[6;80]$
  \task $A[8;51],\; B[9;0],\; C[10;-57]$
  \task $X[-1;0],\; Y[-2;3],\; Z[-3;8]$
  \task $A[0;-63],\; B[9;36],\; C[6;-15]$
  \task $D[1;7],\; E[3{,}5;35{,}75],\; F[-1;-7]$
  \task $G[3;-32],\; H[2;-9],\; I[1;4]$
  \task $J[-2;13],\; K[3;8],\; L[0;5]$
  \task $M[-1;-9],\; N[1;11],\; O[2;27]$
  \task $P[1;3],\; Q[5;15],\; R[-2;15]$
  \task $S[2;7],\; T[-1;4],\; U[0;-1]$
  \task $V[0;-10],\; W[-2;-26],\; X[2;14]$
  \task $K[0;1],\; L[-1;11],\; M[1;11]$
  \task $P\!\left[\frac{5}{2};\frac{229}{4}\right],\; R[3;75],\; S[1;7]$
\end{tasks}

\end{exercise}


% =========================
% ÚLOHA 15
% =========================
\begin{exercise}[Určení kvadratické funkce z podmínek]
Najděte předpis kvadratické funkce \(f\), která splňuje zadané podmínky:

\begin{tasks}

\task $ f(-1) = 0,\quad f(4) = 0,\quad f(0) = 2.$

\task $ f(2) = -3,\quad f(5) = 12,\quad f(-1) = 0.$

\task $ f(1) = 5,\quad f(-2) = 0,\quad f(3) = 14.$

\task $ f(-3)=4,\quad f(0)=1,\quad f(4)=25.$

\end{tasks}

\end{exercise}


% =========================
% BONUSOVÁ ÚLOHA
% =========================

% BONUS 1
\begin{exercise}[Bonus: určení funkce z hodnot a obecného tvaru]
Najděte předpis kvadratické funkce \(f\), která splňuje:

\begin{tasks}
\task $f(1) = 7,\quad f(-2) = -5,\quad f(a) = a(x+1)(x-4)$
\end{tasks}

(Ve třetí podmínce využijte, že po dosazení $x=a$ se obě strany musí rovnat.)
\end{exercise}

% BONUS 2
\begin{exercise}[Bonus: vrchol, hodnota a vztah mezi koeficienty]
Najděte předpis kvadratické funkce \(f\), která splňuje:

\begin{tasks}
\task $V[-1;4],\quad f(2)=3,\quad b = 2a$
\task $V[3;-2],\quad f(0)=10,\quad c = -a$
\end{tasks}
\end{exercise}

% BONUS 3
\begin{exercise}[Bonus: kořen, bod a vrchol]
Určete předpis kvadratické funkce, která splňuje:

\begin{tasks}
\task $f(1)=0,\quad f(5)=7,\quad V[-2;9]$
\task $f(-3)=0,\quad f(2)=13,\quad V[1;-4]$
\end{tasks}
\end{exercise}

% BONUS 4
\begin{exercise}[Bonus: dva body a tvar paraboly]
Najděte kvadratickou funkci, která splňuje:

\begin{tasks}
\task $f(0)=5,\quad f(3)=-4,\quad$ parabola je otevřená dolů
\task $f(-1)=6,\quad f(4)=1,\quad$ osa souměrnosti má rovnici $x=2$
\end{tasks}
\end{exercise}

% BONUS 5
\begin{exercise}[Bonus: různé kombinace podmínek]
Určete funkci \(f\), která splňuje:

\begin{tasks}(2)
\task $f(2)=-3,\quad f(-1)=0,\quad c = 4$
\task $f(4)=20,\quad f(0)=4,\quad$ parabola má maximum
\task $f(-2)=1,\quad f(1)=10,\quad$ vrchol leží na ose $y$
\task $f(3)=7,\quad f(-3)=7,\quad$ graf je sudá funkce
\end{tasks}
\end{exercise}



\end{document}